\documentclass[11pt, oneside]{article} 
\usepackage{geometry}
\geometry{letterpaper} 
\usepackage{graphicx}
	
\usepackage{amssymb}
\usepackage{amsmath}
\usepackage{parskip}
\usepackage{color}
\usepackage{hyperref}

\graphicspath{{/Users/telliott/Dropbox/Github-Math/figures/}}
% \begin{center} \includegraphics [scale=0.4] {gauss3.png} \end{center}

\title{Pythagorean Theorem}
\date{}

\begin{document}
\maketitle
\Large

%[my-super-duper-separator]

\label{sec:pythagorean_thm}

The most famous theorem of Greek geometry is also without a doubt the most useful in calculus.  
\begin{center} \includegraphics [scale=0.2] {pythagoras.png} \end{center}

Pythagoras (c.570 - c.490 BCE) was much younger than Thales but may have encountered him as a youth.  Like many Greek mathematicians, Pythagoras was not from the mainland, but from one of the islands, in his case, Samos.  

Samos is off the coast of what was called Asia Minor (modern Turkey), and Thales lived on the mainland not too far from Samos, at Miletus.  Later, Pythagoras moved to the Italian mainland at Croton.

During his lifetime, Pythagoras was known as a philosopher much more than as a mathematician.  For example, he was famous as an expert on the fate of the soul after death.

\url{https://plato.stanford.edu/entries/pythagoras/}

Wilczek cites Bertrand Russell on Pythagoras:

\begin{quote}"A combination of Einstein and Mary Baker Eddy."\end{quote}

which will only be funny if you look up Mary Baker Eddy.

Pythagoras founded a "school" and it is not sure now which of the theorems developed by this school are due to Pythagoras, and which to his disciples.  It is not even clear whether the Pythagorean theorem, as we understand it today, was known to Pythagoras.  There is no contemporaneous account (Plato, Aristotle) connecting Pythagoras with the theorem.

Regardless of whose idea it was, and who could prove it first, it's clear that they knew something.  Not just the classic 3-4-5 right triangle, but a number of other \emph{Pythagorean triples} had been known for a thousand years.  The tablet "Plimpton 322" contains the triplet 4601-4800-6649 and it dates to about 1800 BCE.

\begin{center} \includegraphics [scale=0.5] {pyth8.png} \end{center}

We're going to spend several chapters examining a few particular proofs of this theorem (there are literally hundreds of them), to look at the ideas from different angles.  This chapter is an introduction that shows some basic algebraic proofs.  

The proof due to the Greeks that we see in the next chapter, is from Euclid.  It employs SAS for triangle congruence, and is probably his own contribution.  Lastly, we'll explore how the concept of area is connected to the Pythagorean theorem.

To begin, let's examine a special case, easily proved, for an isosceles right triangle.

\begin{center} \includegraphics [scale=0.3] {pyth7.png} \end{center}

The area of the square on the hypotenuse is equal to twice the area on each side, in this special case where the sides are equal.

The general theorem is that, that if $a$ and $b$ are the shorter sides of a right triangle, and $c$ is the hypotenuse
\[ a^2 + b^2 = c^2 \]

The following one is sometimes called the "Chinese proof."  I can easily imagine proceeding from the figure above to the left panel below by simply rotating the inner square and collapsing the surrounding one.

\begin{center} \includegraphics [scale=0.35] {pythagoras1.png} \end{center}

It really needs no explanation, but ..

We have a large square box that contains within it a white square, whose side is also the hypotenuse of the four identical right triangles contained inside.  Altogether the four triangles plus the white area add up to the total.

We simply rearrange the triangles.  Now we evidently have the same area left over from the four triangles, because they still have the same area and the surrounding box has not changed.  

But clearly, now the white area is the sum of the squares on the second and third sides of the triangles.  Hence the two white squares on the right are equal in area to the large white square on the left.  

$\square$

This diagram is contained in the Chinese text Zhoubi Suanjing.

\begin{center} \includegraphics [scale=0.25] {Chinese_pythagoras.jpg} \end{center}

Eight right triangles are formed with sides of $3$ and $4$ units.  There are $7$ units on each side of the large square so the total area is $49$.  Each pair of triangles has area $12$ ($3 \cdot 4$), so the square in the center really is a unit square with area $49 - 4 \cdot 12 = 1$.

So then, the central square consists of four triangles with total area $24$ plus the unit square for $25$.  This is the square of the long side, namely $5$.

\url{https://en.wikipedia.org/wiki/Zhoubi_Suanjing}

There is debate about whether this is really a proof, or if it simply presents the arithmetic for the example of a 3-4-5 right triangle.

\subsection*{simple proofs}

Many proofs of the theorem are algebraic.  Here are two:

\begin{center} \includegraphics [scale=0.25] {pythagoras8.png} \end{center}

\emph{Proof}.

In the left panel, we have arranged four identical right triangles.  The white square at the center has sides of length $c$ and angles that are right angles because when summed to two complementary angles, the result is two right angles.

The algebra is
\[ (a + b)^2 = 4 \cdot \frac{1}{2} \cdot ab + c^2 \]
\[ a^2 + 2ab + b^2 = 2 ab + c^2 \]

Subtract $2ab$ from both sides, and we're done.  

$\square$

I will leave the right panel to you.  It is pretty easy.

\subsection*{similar triangles}

\label{sec:Pythagoras_similar_triangles}

\label{sec:geometric_mean_pyth}

Here is a simple classic that depends on the ratios formed for similar right triangles:

\begin{center} \includegraphics [scale=0.4] {triangle3.png} \end{center}

\emph{Proof}

We know that when an altitude is drawn in a right triangle, the two resulting right triangles are similar, by complementary angles.  Similarity means that we have equal ratios of sides.  Here are two sets:

hypotenuse to short side
\[ \frac{a}{x} = \frac{b}{h} = \frac{c}{a} \]
hypotenuse to long side
\[ \frac{a}{h} = \frac{b}{y} = \frac{c}{b} \]

From the first
\[ a^2 = cx \]
And from the second
\[ b^2 = cy \]

Add them:

\[ a^2 + b^2 = cx + cy \]
\[ = c(x+y) = c^2 \]

$\square$

Another relationship from similar triangles is long side to short side:
\[ \frac{h}{x} = \frac{y}{h} \]
\[ h^2 = xy \]
\[ h = \sqrt{xy} \]

$h$ is the \emph{geometric mean} of $x$ and $y$.

\subsection*{converse of Pythagorean theorem}

In reasoning deductively, we move from premises (data given, previous theorems that were proved), and use logic to reach a conclusion.  The question arises whether, if the conclusion is given, does it follow that the premises are true?  This is the problem of the \emph{converse} of a theorem.

It may be so, or it may not.

We can state the converse of the Pythagorean theorem as follows:  suppose we have triangle such that $a^2 + b^2 = c^2$.  Does it follow that the angle between $a$ and $b$ is a right angle?

We profess to not know whether it is or is not a right angle.

\emph{Proof}.

So then, draw another triangle with sides $a$ and $b$ that \emph{is} a right triangle.  Then, by the forward theorem, $a^2 + b^2 = c^2$.  

But the second triangle is congruent to the first one by SSS, each side is the same.  Since they are equivalent parts of congruent triangles, the angle between $a$ and $b$ is a right angle. 

$\square$

\subsection*{David Bowie problem}

Here's a problem from the web:

\begin{center} \includegraphics [scale=0.4] {bowie1.png} \end{center}

One way to look at this is to imagine the octagon inscribed in a circle of diameter $d = 2r$ (below).  We reason that the two triangles are right triangles with a shared hypotenuse.  

\begin{center} \includegraphics [scale=0.4] {bowie2.png} \end{center}

We can apply trigonometry to find that the whole area required is simply
\[ A_{\text{ region}} = d^2 \sin \phi \cos \phi \]

But we aren't ready to do this kind of calculation yet.

However, it will turn out that the part shaded red is simply one-half of the whole.  

Therefore, the three lines divide the octagon into four equal parts.  That simple answer is a clue that something important makes the problem easy.

\begin{center} \includegraphics [scale=0.4] {bowie3.png} \end{center}

The triangle with small sides $a$ can be moved so one of the unshaded base parts becomes a rectangle with area
\[ a(s + a) = sa + a^2 \]

while the base of the triangle is $s + 2a$ so its area is
\[ \frac{1}{2} s (s + 2a) = \frac{1}{2} s^2 + sa \]

It's not obvious at first that these are equal, but working with it, they would be equal if we can show that $2a^2 = s^2$.  

Of course we can do exactly that, because $a$ is the side of an isosceles right triangle with hypotenuse $s$.  By the Pythagorean theorem, we have
\[ a^2 + a^2 = s^2 \]

$\square$
 
\end{document}